\section{Introduction}
\begin{figure*}[t]
  \centering
  \includegraphics[width=\textwidth]{headline_echo.pdf}
  \caption{The echo and its removal. (a)~In synchronous SCA, the echo of a spin's state before a flip pushes the spin back. (b)~Anechoic subtracts the echo from every decision. (c)~Share of flips that undo the previous step on K2000 (engine model, 1{,}024 runs per schedule), with TTS measured on the V80 (Table~\ref{tab:configs}).}
  \Description{Three panels. Panel a shows a spin that flips between steps t minus one and t; its neighbors respond to its state before the flip, and their response returns in its local field and flips the spin back. Panel b shows the same sequence with the returning term subtracted, so the flip is kept. Panel c plots the percentage of flips that undo the previous step against the annealing step: 67 percent on average for plain SCA with STATICA's 560-step schedule, whose time-to-solution on the V80 is 0.231 ms, and 37 percent for Anechoic with 280 steps, whose time-to-solution is 0.050 ms.}
  \label{fig:headline}
\end{figure*}

Ising machines solve combinatorial optimization problems by letting a spin system relax toward low energy. Optical, analog, quantum and digital designs have all been demonstrated~\cite{doi:10.1126/science.aah4243,goto2019combinatorial,9162869,boothby2020next,9062965,10609617}. Dense problems are particularly demanding for hardware: every spin interacts with every other spin, so the problem cannot be embedded into local connectivity, and every update affects every local field. The standard benchmark for such problems is K2000, Max-Cut on a complete graph with 2{,}000 spins~\cite{doi:10.1126/science.aah4243}. Machines are compared by their time-to-solution (TTS), the time required to reach a target cut with 99\% probability. The fastest published hardware result is 0.26\,ms, from Toshiba's ballistic simulated-bifurcation machine (bSB) on an FPGA~\cite{goto2021high}.

Three families of dynamics trade work per step against progress per step: (1)~\emph{simulated annealing} (SA)~\cite{kirkpatrick1983optimization}, which updates one spin at a time and descends reliably, although on a dense graph each update rewrites the local fields of all $N$ spins, so a sweep is serial; (2)~\emph{simulated bifurcation}~\cite{goto2019combinatorial,goto2021high}, whose ballistic form (bSB) needs the fewest steps, although each step is a real-valued matrix--vector product, dense until the amplitudes settle; and (3)~\emph{synchronous stochastic cellular automata} (SCA)~\cite{9062965}, which decide all spins in parallel and access the couplings only for the spins that flip, giving the least expensive step in hardware. The limitation of SCA lies in its dynamics: a spin that flips is pushed back by the echo of its own previous state through the network, and on K2000 two-thirds of all flips reverse the flip of the previous step (Figure~\ref{fig:headline}).

This paper presents Anechoic, an FPGA Ising machine that combines this inexpensive step with a correction that subtracts the echo, as an anechoic chamber absorbs reflections (Figure~\ref{fig:headline}b). Dynamical mean-field theory~\cite{sompolinsky1982relaxational,eissfeller1992new,coolen2001statistical} identifies the echo as an Onsager reaction term, whose lag-one part has the coefficient $n_{\mathrm{lin}}/2T$, where $n_{\mathrm{lin}}$ is the number of spins whose flip probability lies strictly between 0 and 1 and $T$ is the temperature. The decision logic already flags these spins, so subtracting this lag-one term costs one population count per step. As an SCA step's cost grows with its flips, the correction saves time twice: the anneal needs fewer steps and starts colder, with fewer flips.

The contributions of this work are:
\begin{itemize}[leftmargin=*]
\item \textbf{Onsager-corrected SCA} (Section~\ref{sec:algorithm}). The echo is subtracted with a counted coefficient (Onsager-online) or a temperature-scaled one (Onsager-$\kappa T$); at every step budget on K2000, both forms reach lower Ising energy than five published binary-spin algorithms run at their papers' settings.
\item \textbf{A flip-proportional dense engine} (Section~\ref{sec:arch}). A hand-written RTL engine for the AMD Alveo V80~\cite{amd_v80} executes a dense 2{,}000-spin step in \cycStep{} cycles plus \cycFlip{} cycles per flip, bit-exact with a reference model in simulation and on the board.
\item \textbf{A measured 12-engine system} (Section~\ref{sec:evaluation}). Twelve engines reach a cut of at least 33{,}000 on K2000 with 99\% probability in \ttsP{}, \spdBSB{} faster than the fastest published hardware~\cite{goto2021high}, at \ets{} per solution. With 2-bit couplings ($-1$, 0, $+1$), Onsager-$\kappa T$ is also faster than plain SCA on 48 of 51 G-set instances.
\item \textbf{Cross-platform evidence} (Sections~\ref{subsec:gpu} and~\ref{subsec:asic}). On an NVIDIA GH200 GPU~\cite{nvidia_gh200}, the bit-exact engine is \gpuSpd{} slower and needs \gpuEratio{} more energy per solution, though its 240 concurrent anneals give higher throughput. As a 22\,nm ASIC routed at 1.25\,GHz, twelve engines would reach the target in about 10\,\textmu s, $26\times$ faster than the fastest published hardware.
\end{itemize}
